
\section{Enquadramento e objetivos}
\label{sec:enq-e-obj}

Os manipuladores robóticos desempenham um papel importante na automação industrial, sendo utilizados em tarefas que exigem precisão, repetibilidade e flexibilidade de movimentos. O seu funcionamento depende do conhecimento da geometria do manipulador e da relação entre as variáveis articulares e o posicionamento do \textit{end effector} no espaço de trabalho.

Neste contexto, o presente projeto centra-se no estudo do manipulador myCobot~280, recorrendo ao MATLAB para implementar os algoritmos cinemáticos e ao CoppeliaSim para simular e avaliar o comportamento do robô.

Os principais objetivos do trabalho são:

\begin{itemize}
    \item Desenvolver a formulação matemática da cinemática direta e inversa do manipulador;

    \item Implementar os modelos cinemáticos em MATLAB;

    \item Determinar a posição do centro de preensão e a direção da garra a partir dos ângulos articulares;

    \item Calcular as configurações articulares necessárias para alcançar posições pretendidas, considerando a inclinação da garra e os limites das juntas;

    \item Implementar movimentos no espaço das juntas e trajetórias no espaço cartesiano;

    \item Integrar o MATLAB com o CoppeliaSim para executar e visualizar os movimentos do manipulador;

    \item Explorar a utilização da visão artificial na identificação e localização de objetos;

    \item Validar os modelos implementados através da comparação entre os resultados calculados e os obtidos em simulação.
\end{itemize}
